\scanpage{20}
\mspage{11}
% archive sheet number: 20
% lines 1--4
Les \deleted{structures}\inserted{relations} $\leq$, $\ll$, $\unsure{\gtrless|\bowtie}$, $\parallel$ sur $\unsure{\mathscr{F}|\mathcal{F}}_{M^{\unsure{d|a|\circ}}}$ induisent des \deleted{$\unsure{\Rightarrow}$}\deleted{structures}\inserted{relations} de même nom sur $\mathscr{F}$.
Les prop $\unsure{1_{\circ}\mathrm{F}|1^{\circ}\mathrm{F}|1_0\mathrm{F}|1_o\mathrm{F}}$, $2_{\circ}\mathrm{F}$ (p. 2 bis), prop. 2, 3, $\deleted{\unsure{3}}5$ \unsure{pour|passent} ``ipso facto'' \unsure{se conservent} dans un atlas général.

% lines 5--8
Il n'y a que la prop. 4 (recollement des raffinements) qui, dans le cas d'un ens d'indices ``\uline{infini}'' $I$, demande un axiome de stabilité supplémentaire pour l'atlas,

% lines 9--14
Savoir\par
\uline{At Fil}\quad Soit $F\in\mathscr{F}$, et soit $(\deleted{\unsure{F|E}}G_\alpha)$ une famille filtrante croissante \inserted{pour $\leq$} de raffinements de $F$. Supposons que $\forall X\in\widetilde{F}$, la famille des $G_\alpha\mid X$ soit stationnaire. Alors \deleted{$\unsure{\operatorname{Sup}G_\alpha|\bigcup G_\alpha}$} $\operatorname{Sup}^{\leq}G_\alpha$ existe dans $\mathscr{F}$ \deleted{(et est un raffinement} \deleted{\unsure{(i.e.}}.

% lines 15--18
Du point de vue $\mathscr{F}\subset\mathscr{F}_M$, cela signifie que $\bigcup G_\alpha$ (qui est le Sup dans $\mathscr{F}_M$) est en fait dans $\mathscr{F}$ : propriété de stabilité de $\mathscr{F}$ en \deleted{\unsure{tant}} que partie fermée de $\mathscr{F}_M,\leq$.

% lines 19--22
On peut considérer la structure d'atlas comme donnée indépendamment de celle d'un magasin donné au préalable.

Atlas : $\mathscr{F}$ muni de trois relations $\leq$, $\ll$, $\unsure{\mid\circ\mid|\parallel}$

% lines 23--27
\uline{At 1}\quad Les Sup majorés existent, et $\mathscr{F}\neq\varnothing$ ($\Rightarrow\mathscr{F}$ a un plus petit élément $\varnothing_{\mathscr{F}}$).

\deleted{\unsure{Éléments}} Figures élémentaires \unsure{$\cup$|$\vee$}-éléments strictement irréductibles :
\[
X=\operatorname{Sup}F_i\Rightarrow\exists i,\ \deleted{\unsure{K_i}}F_i=X.
\]
$\deleted{\unsure{E}}M$ ens des fig. élém., $\widetilde{F}=\{X\in M\mid X\leq F\}$.

% slanted
\margin{At 2, 3, 4 : \uline{passage aux strates}}
% lines 28--30
\[
\text{\uline{At 2}}\quad F\leq G\Longleftrightarrow\widetilde{\deleted{\unsure{f}}F}\subset\widetilde{G}\qquad\deleted{\unsure{(i.e.)}}\text{i.e. }X\in M,\ X\leq F\Rightarrow X\leq G.
\]
``toute strate de $F$ est \unsure{strate} de $G$''.

\unsure{Equivalent|Équivalent} : Toute figure est le Sup de ses strates
